\begin{tabularx}{\textwidth}[t]{|X|X|}
  \hline 
  \textbf{Notions} & \textbf{Commentaires} \\
  \hline 
  Algorithme des $k$ plus proches voisins avec distance euclidienne.
   & 
   Matrice de confusion. Lien avec l'apprentissage supervisé.
  \\
  \hline 
  Algorithme des  k-moyennes.
   & 
   Lien avec l'apprentissage non-supervisé.
La démonstration de la convergence n'est pas au programme. On observe des convergences vers des minima locaux.
  \\
  \hline 
  \multicolumn{2}{|X|}{\textbf{Mise en œuvre}}
  \\
  \hline 
  \multicolumn{2}{|p{0.95\textwidth}|}{
  La connaissance dans le détail des algorithmes de cette section n'est pas un attendu du programme. Les étudiants acquièrent une familiarité avec les idées sous-jacentes qu'ils peuvent réinvestir dans des situations où les modélisations et les recommandations d'implémentation sont guidées, notamment dans leurs aspects arborescents.
  } 
  \\
  \hline
\end{tabularx}
